\chapter{Mapping Cones}
\label{ch:viii07-mapping-cones}

Cell attachment is a special case of a more flexible construction.  Given a
map \(f:A\to X\), instead of attaching only a disk along a sphere, one can
attach the whole cone on \(A\).  The resulting mapping cone records the map
itself as part of a new space.  Mapping cones organize quotient spaces,
cofibers, suspensions, and many exact-sequence constructions that appear later
in homology.

\section{The cone on a space}

\begin{definition}[Cone]
\label{def:viii07-cone}
For a space \(A\), the cone on \(A\) is
\[
CA=(A\times I)/(A\times\{1\}),
\]
where all points of \(A\times\{1\}\) are identified to one cone vertex.
\end{definition}

The copy \(A\times\{0\}\) is the base of the cone.

\begin{proposition}[Cones are contractible]
\label{prop:viii07-cone-contractible}
For every nonempty space \(A\), the cone \(CA\) is contractible.
\end{proposition}

\begin{proof}
Push every point linearly toward the cone vertex in the \(I\)-coordinate.
The quotient description makes the terminal slice a single point.
\end{proof}

\section{The suspension}

\begin{definition}[Suspension]
\label{def:viii07-suspension}
The suspension of \(A\) is
\[
\Sigma A=(A\times I)/(A\times\{0\},A\times\{1\}),
\]
where each end is collapsed to a point.
\end{definition}

Equivalently, \(\Sigma A\) is obtained by gluing two cones on \(A\) along their
common base.  In particular,
\[
\Sigma S^n\cong S^{n+1}.
\]

\section{Definition of the mapping cone}

\begin{definition}[Mapping cone]
\label{def:viii07-mapping-cone}
For a continuous map \(f:A\to X\), the mapping cone is
\[
C_f=X\cup_f CA,
\]
where the base \(A\subset CA\) is identified with \(X\) using \(f\).
\end{definition}

Thus \(C_f\) is the pushout
\[
\begin{array}{ccc}
A&\xrightarrow{f}&X\\
\big\downarrow&&\big\downarrow\\
CA&\longrightarrow&C_f.
\end{array}
\]

\section{Cell attachment as a mapping cone}

If \(A=S^{n-1}\), then
\[
CA\cong D^n.
\]
Therefore
\[
X\cup_\varphi D^n
\cong
C_\varphi
\]
for every attaching map \(\varphi:S^{n-1}\to X\).

\begin{proposition}[Cell-attachment specialization]
\label{prop:viii07-cell-special}
Ordinary \(n\)-cell attachment is the mapping-cone construction for a map
whose domain is \(S^{n-1}\).
\end{proposition}

\section{Mapping cone of the identity}

\begin{proposition}[Cone of the identity]
\label{prop:viii07-id}
For every space \(A\),
\[
C_{\operatorname{id}_A}\cong CA.
\]
Hence it is contractible.
\end{proposition}

The copy of \(A\) already present as \(X=A\) is exactly the base to which the
cone is attached.

\section{Mapping cone of a constant map}

\begin{proposition}[Constant map cone]
\label{prop:viii07-constant}
If \(f:A\to X\) is constant at a chosen basepoint \(x_0\), then under standard
pointed hypotheses,
\[
C_f\simeq X\vee\Sigma A.
\]
In the literal quotient model, the cone base is attached at one point of \(X\).
\end{proposition}

For \(X=\{*\}\),
\[
C_f\cong \Sigma A.
\]

\section{Mapping cone of an inclusion}

Let \(i:A\hookrightarrow X\) be an inclusion.

\begin{proposition}[Cone of a cofibration-like inclusion]
\label{prop:viii07-inclusion}
When \(A\subset X\) is a CW subcomplex, the mapping cone \(C_i\) is naturally
homotopy equivalent to the quotient
\[
X/A.
\]
\end{proposition}

For CW pairs this is one of the main reasons mapping cones model relative
topology.

\section{The cofiber viewpoint}

The sequence
\[
A\xrightarrow{f}X\longrightarrow C_f
\]
is called a cofiber sequence at its first stage.  After suspending, one obtains
the extended pattern
\[
A\xrightarrow{f}X\longrightarrow C_f\longrightarrow\Sigma A
\longrightarrow\Sigma X\longrightarrow\cdots.
\]
Later homology turns this geometric sequence into exact algebraic sequences.

\section{Examples from spheres}

For the degree-\(m\) map
\[
f_m:S^1\to S^1,\qquad f_m(z)=z^m,
\]
the mapping cone is
\[
C_{f_m}=S^1\cup_{f_m}D^2.
\]
Thus:
\[
C_{f_0}\simeq S^1\vee S^2,\qquad
C_{f_1}\cong D^2,\qquad
C_{f_2}\cong\mathbb RP^2.
\]

\section{Moore-space preview}

A basic Moore space can be obtained by attaching an \((n+1)\)-cell to \(S^n\)
along a degree-\(m\) map:
\[
M(\mathbb Z/m,n)
=
S^n\cup_{m}D^{n+1}.
\]
Its cellular homology later realizes a single \(\mathbb Z/m\) group in the
desired degree.

\section{Functorial behavior}

Suppose there is a commutative square
\[
\begin{array}{ccc}
A&\xrightarrow{f}&X\\
\downarrow u&&\downarrow v\\
B&\xrightarrow{g}&Y.
\end{array}
\]
The maps \(u\) and \(v\) induce a natural map
\[
C_f\to C_g
\]
by applying \(v\) on the old space and the cone on \(u\) on the cone part.

\begin{proposition}[Map of mapping cones]
\label{prop:viii07-functorial}
A strictly commutative square of maps induces a continuous map between the
corresponding mapping cones.
\end{proposition}

\section{Quotient description}

One can write
\[
C_f=
(X\sqcup A\times I)/\!\sim,
\]
where
\[
(a,0)\sim f(a)
\]
and all \((a,1)\) are identified to the cone vertex.

This description is useful for checking maps explicitly.

\section{Mapping cylinders versus mapping cones}

\begin{definition}[Mapping cylinder]
\label{def:viii07-mapping-cylinder}
The mapping cylinder of \(f:A\to X\) is
\[
M_f=(A\times I)\sqcup X/\bigl((a,0)\sim f(a)\bigr).
\]
\end{definition}

The mapping cone is obtained from \(M_f\) by collapsing the free end
\(A\times\{1\}\) to a point.

\begin{proposition}[Cylinder retract]
\label{prop:viii07-cylinder-retract}
The mapping cylinder \(M_f\) deformation retracts onto the copy of \(X\).
\end{proposition}

\section{Null-homotopic maps}

If \(f:A\to X\) is null-homotopic, the mapping cone is closely related to
\(X\vee\Sigma A\).  For CW-type spaces this relationship is a homotopy
equivalence.

\begin{theorem}[Null-map splitting]
\label{thm:viii07-null-splitting}
For a based null-homotopic map between suitable CW spaces,
\[
C_f\simeq X\vee\Sigma A.
\]
\end{theorem}

\section{Homotopy invariance preview}

If \(f_0\simeq f_1:A\to X\), one expects
\[
C_{f_0}\simeq C_{f_1}.
\]
This is the mapping-cone version of the homotopic-attaching-map theorem.
VIII/08 proves the relevant principle in the cell-attachment setting.

\begin{warning}[Not automatically homeomorphic]
\label{warn:viii07-homeomorphic}
Homotopic maps need not produce literally identical quotient spaces.  The
natural invariant conclusion is homotopy equivalence, not homeomorphism.
\end{warning}


\section{Exercises}

\begin{exercise}\label{exr:viii07-01}
Define the cone \(CA\).
\end{exercise}
\begin{hint}
Collapse one end of \(A\times I\).
\end{hint}
\begin{solution}
\(CA=(A\times I)/(A\times\{1\})\).
\end{solution}

\begin{exercise}\label{exr:viii07-02}
Why is \(CA\) contractible?
\end{exercise}
\begin{hint}
Move every point toward the cone vertex.
\end{hint}
\begin{solution}
The homotopy that increases the \(I\)-coordinate to \(1\) contracts the cone to its vertex.
\end{solution}

\begin{exercise}\label{exr:viii07-03}
Define \(\Sigma A\).
\end{exercise}
\begin{hint}
Collapse both ends.
\end{hint}
\begin{solution}
\(\Sigma A=(A\times I)/(A\times\{0\},A\times\{1\})\).
\end{solution}

\begin{exercise}\label{exr:viii07-04}
Express \(S^{n+1}\) as a suspension.
\end{exercise}
\begin{hint}
Suspend \(S^n\).
\end{hint}
\begin{solution}
\(S^{n+1}\cong\Sigma S^n\).
\end{solution}

\begin{exercise}\label{exr:viii07-05}
Define \(C_f\) for \(f:A\to X\).
\end{exercise}
\begin{hint}
Attach the cone base using \(f\).
\end{hint}
\begin{solution}
\(C_f=X\cup_f CA\).
\end{solution}

\begin{exercise}\label{exr:viii07-06}
How is an \(n\)-cell attachment a mapping cone?
\end{exercise}
\begin{hint}
Use \(CS^{n-1}\cong D^n\).
\end{hint}
\begin{solution}
\(X\cup_\varphi D^n\cong C_\varphi\).
\end{solution}

\begin{exercise}\label{exr:viii07-07}
What is \(C_{\operatorname{id}_A}\)?
\end{exercise}
\begin{hint}
The cone is attached to its own base copy.
\end{hint}
\begin{solution}
It is homeomorphic to \(CA\), hence contractible.
\end{solution}

\begin{exercise}\label{exr:viii07-08}
What is the mapping cone of a constant map to a point?
\end{exercise}
\begin{hint}
The target is one point.
\end{hint}
\begin{solution}
It is \(\Sigma A\).
\end{solution}

\begin{exercise}\label{exr:viii07-09}
For a CW subcomplex inclusion \(A\hookrightarrow X\), what quotient models the cone?
\end{exercise}
\begin{hint}
Collapse \(A\).
\end{hint}
\begin{solution}
\(C_i\simeq X/A\).
\end{solution}

\begin{exercise}\label{exr:viii07-10}
Write the first three terms of a cofiber sequence.
\end{exercise}
\begin{hint}
Start with \(f\).
\end{hint}
\begin{solution}
\(A\xrightarrow f X\to C_f\).
\end{solution}

\begin{exercise}\label{exr:viii07-11}
What follows \(C_f\) in the extended cofiber sequence?
\end{exercise}
\begin{hint}
Suspension of the domain.
\end{hint}
\begin{solution}
\(\Sigma A\).
\end{solution}

\begin{exercise}\label{exr:viii07-12}
Identify \(C_{z\mapsto z}\) for \(S^1\to S^1\).
\end{exercise}
\begin{hint}
Degree one attachment.
\end{hint}
\begin{solution}
It is \(D^2\).
\end{solution}

\begin{exercise}\label{exr:viii07-13}
Identify \(C_{z\mapsto z^2}\).
\end{exercise}
\begin{hint}
Recall the projective-plane attachment.
\end{hint}
\begin{solution}
It is \(\mathbb RP^2\).
\end{solution}

\begin{exercise}\label{exr:viii07-14}
Identify the constant-map cone \(S^1\to S^1\) up to homotopy.
\end{exercise}
\begin{hint}
Use the null-map splitting.
\end{hint}
\begin{solution}
\(S^1\vee S^2\).
\end{solution}

\begin{exercise}\label{exr:viii07-15}
What is a Moore-space attachment?
\end{exercise}
\begin{hint}
Attach by a degree-\(m\) map.
\end{hint}
\begin{solution}
\(S^n\cup_mD^{n+1}\).
\end{solution}

\begin{exercise}\label{exr:viii07-16}
Define the mapping cylinder \(M_f\).
\end{exercise}
\begin{hint}
Do not collapse the free end.
\end{hint}
\begin{solution}
\((A\times I)\sqcup X/((a,0)\sim f(a))\).
\end{solution}

\begin{exercise}\label{exr:viii07-17}
How do you get \(C_f\) from \(M_f\)?
\end{exercise}
\begin{hint}
Collapse the free end.
\end{hint}
\begin{solution}
Collapse \(A\times\{1\}\) to one point.
\end{solution}

\begin{exercise}\label{exr:viii07-18}
To what does \(M_f\) deformation retract?
\end{exercise}
\begin{hint}
Push the cylinder into the glued end.
\end{hint}
\begin{solution}
To the copy of \(X\).
\end{solution}

\begin{exercise}\label{exr:viii07-19}
What does a commutative square induce on cones?
\end{exercise}
\begin{hint}
Map the target part and cone part separately.
\end{hint}
\begin{solution}
A natural map \(C_f\to C_g\).
\end{solution}

\begin{exercise}\label{exr:viii07-20}
What splitting holds for a suitable null-homotopic map?
\end{exercise}
\begin{hint}
Use suspension.
\end{hint}
\begin{solution}
\(C_f\simeq X\vee\Sigma A\).
\end{solution}

\begin{exercise}\label{exr:viii07-21}
Are mapping cones of homotopic maps necessarily homeomorphic?
\end{exercise}
\begin{hint}
The expected invariant is weaker.
\end{hint}
\begin{solution}
No; the natural conclusion is homotopy equivalence.
\end{solution}

\begin{exercise}\label{exr:viii07-22}
Write the quotient model of \(C_f\).
\end{exercise}
\begin{hint}
Use \(X\sqcup A\times I\).
\end{hint}
\begin{solution}
Identify \((a,0)\) with \(f(a)\) and collapse all \((a,1)\) to the cone vertex.
\end{solution}

\begin{exercise}\label{exr:viii07-23}
Why are mapping cones useful for relative topology?
\end{exercise}
\begin{hint}
Consider a subcomplex inclusion.
\end{hint}
\begin{solution}
For \(A\subset X\), \(C_i\simeq X/A\), so relative information becomes an ordinary quotient space.
\end{solution}

\begin{exercise}\label{exr:viii07-24}
What is the bridge from VIII/07 to VIII/08?
\end{exercise}
\begin{hint}
Vary the attaching map through a homotopy.
\end{hint}
\begin{solution}
VIII/08 proves that homotopic attaching maps give the same homotopy type under the standard CW hypotheses.
\end{solution}


\section{Solved problem dossiers}

\begin{problem}[Cone contraction]\label{prob:viii07-01}
Construct an explicit contraction of \(CA\) to its vertex.
\end{problem}
\begin{solution}
Represent a cone point by \([a,s]\).  Define \(H([a,s],t)=[a,(1-t)s+t]\).  At \(t=1\) every point has coordinate \(1\), hence is the cone vertex.
\end{solution}

\begin{problem}[Suspension of a point]\label{prob:viii07-02}
Identify \(\Sigma\{*\}\).
\end{problem}
\begin{solution}
The interval has each endpoint separately collapsed, which changes nothing; the result is an interval, hence contractible.
\end{solution}

\begin{problem}[Suspension of \(S^0\)]\label{prob:viii07-03}
Show \(\Sigma S^0\cong S^1\).
\end{problem}
\begin{solution}
Two intervals, one for each point of \(S^0\), share both suspension endpoints.  Their union is a circle.
\end{solution}

\begin{problem}[Identity cone]\label{prob:viii07-04}
Prove \(C_{\operatorname{id}_A}\cong CA\).
\end{problem}
\begin{solution}
The target copy \(A\) is glued pointwise to the cone base.  It contributes no additional points beyond that base, leaving precisely the cone.
\end{solution}

\begin{problem}[Constant cone]\label{prob:viii07-05}
Compute the mapping cone of the constant map \(S^{n-1}\to\{*\}\).
\end{problem}
\begin{solution}
It is the suspension \(\Sigma S^{n-1}\cong S^n\).
\end{solution}

\begin{problem}[Degree-one circle cone]\label{prob:viii07-06}
Compute \(C_f\) for \(f:S^1\to S^1\) of degree \(1\) represented by the identity.
\end{problem}
\begin{solution}
The cone is a disk attached along its full boundary by the identity, so the result is \(D^2\), hence contractible.
\end{solution}

\begin{problem}[Degree-two circle cone]\label{prob:viii07-07}
Explain why the cone of \(z\mapsto z^2\) is \(\mathbb RP^2\).
\end{problem}
\begin{solution}
The cone on \(S^1\) is \(D^2\); attaching its boundary to \(S^1\) by the degree-two quotient map is the standard cell model of \(\mathbb RP^2\).
\end{solution}

\begin{problem}[Null circle cone]\label{prob:viii07-08}
Compute the cone of a constant \(S^1\to S^1\).
\end{problem}
\begin{solution}
The disk boundary collapses at one point of the old circle, producing \(S^1\vee S^2\).
\end{solution}

\begin{problem}[Inclusion cone]\label{prob:viii07-09}
For the inclusion \(S^{n-1}\hookrightarrow D^n\), identify the quotient model \(D^n/S^{n-1}\).
\end{problem}
\begin{solution}
It is \(S^n\).  Thus the mapping cone of the boundary inclusion has homotopy type \(S^n\).
\end{solution}

\begin{problem}[Pair quotient]\label{prob:viii07-10}
Let \(A\) be a CW subcomplex of \(X\).  Explain why \(C_i\) models \(X/A\).
\end{problem}
\begin{solution}
The cone on \(A\) contracts the entire attached copy of \(A\) to the cone vertex while the base is identified with the original \(A\subset X\).  For a cofibration such as a CW-subcomplex inclusion, this yields the quotient homotopy type.
\end{solution}

\begin{problem}[Mapping-cylinder retraction]\label{prob:viii07-11}
Write an explicit deformation retraction \(M_f\to X\).
\end{problem}
\begin{solution}
On the cylinder send \([a,s]\) to \([a,(1-t)s]\) during the homotopy, keeping \(X\) fixed.  At \(t=1\) every cylinder point lies at \(s=0\), identified with \(f(a)\in X\).
\end{solution}

\begin{problem}[Cone from cylinder]\label{prob:viii07-12}
Describe \(C_f\) as a quotient of \(M_f\).
\end{problem}
\begin{solution}
Collapse the free copy \(A\times\{1\}\subset M_f\) to a single point.  The result is exactly the mapping cone.
\end{solution}

\begin{problem}[Map of cones]\label{prob:viii07-13}
Given \(vf=gu\), construct \(C_f\to C_g\).
\end{problem}
\begin{solution}
Use \(v\) on \(X\) and send \([a,t]\in CA\) to \([u(a),t]\in CB\).  Commutativity ensures agreement on the glued base.
\end{solution}

\begin{problem}[Moore-space preview]\label{prob:viii07-14}
For \(M=S^n\cup_mD^{n+1}\), predict the nontrivial reduced homology group.
\end{problem}
\begin{solution}
Cellular homology later gives a boundary map multiplication by \(m\), so the reduced homology in degree \(n\) is \(\mathbb Z/m\).
\end{solution}

\begin{problem}[Cofiber continuation]\label{prob:viii07-15}
Explain geometrically why a map \(C_f\to\Sigma A\) exists.
\end{problem}
\begin{solution}
Collapse the copy of \(X\subset C_f\) to a point.  The cone part then has its base collapsed as well as its vertex, producing the suspension \(\Sigma A\).
\end{solution}

\begin{problem}[Null-map splitting]\label{prob:viii07-16}
Why should a null-homotopy of \(f:A\to X\) lead to \(C_f\simeq X\vee\Sigma A\)?
\end{problem}
\begin{solution}
The null-homotopy lets the cone base attachment deform to a constant attachment.  A constant attachment wedges the suspended domain onto \(X\).
\end{solution}

\begin{problem}[Two homotopic maps]\label{prob:viii07-17}
Suppose \(f_0\simeq f_1:A\to X\).  What mapping-cone relation should be expected?
\end{problem}
\begin{solution}
Under the standard CW/cofibration hypotheses, \(C_{f_0}\simeq C_{f_1}\).  VIII/08 develops the corresponding attaching-map theorem.
\end{solution}

\begin{problem}[Quotient map check]\label{prob:viii07-18}
Verify that the natural \(X\to C_f\) is continuous.
\end{problem}
\begin{solution}
It is the composite of the inclusion \(X\hookrightarrow X\sqcup CA\) with the quotient map defining \(C_f\), hence continuous.
\end{solution}

\begin{problem}[Suspension as double cone]\label{prob:viii07-19}
Show that \(\Sigma A\) is obtained from two copies of \(CA\) glued along \(A\).
\end{problem}
\begin{solution}
Split \(A\times I\) at \(t=1/2\).  Each half becomes a cone after collapsing its outer end, and their common middle slice is a copy of \(A\).
\end{solution}

\begin{problem}[Bridge to exact sequences]\label{prob:viii07-20}
Why does the cofiber sequence foreshadow homology exact sequences?
\end{problem}
\begin{solution}
Homology converts the geometric passage \(A\to X\to C_f\to\Sigma A\) into groups linked by induced maps; suspension shifts degree, yielding the connecting morphisms of the long exact sequence.
\end{solution}
